\section*{TL;DR}
\addcontentsline{toc}{section}{TL;DR}

\pqhash{} is a 73-byte, lexicographically-sortable \textbf{content identifier} with a 512-bit SHAKE-256 \citep{fips202} hash tail and 9 bytes of provenance metadata (48-bit millisecond timestamp, variant byte, 16-bit counter). It is \emph{not} a hash function, \emph{not} a MAC, \emph{not} a signature, and \emph{not} an encryption envelope. It is a \emph{content identifier} that carries hash-derived collision resistance plus enough metadata to make the identifier itself sortable and traceable to a generator and time.

\pqhash{} is the second of three projects in the \pqtoolkit{} monorepo. It shares its 9-byte header layout with \pqid{} \citep{rfc4648,grover1996} (41-byte identifier, variants \code{0x01}/\code{0x02}) and \pqcrypt{} (ML-KEM-1024 \citep{fips203} encryption envelope, variant \code{0x04}); \pqhash{} owns variant byte \code{0x03}. Cross-project decoders reject every byte they do not own, by KAT contract.

\textbf{Wire layout (73 bytes total):}
\begin{lstlisting}[language=Pseudocode]
Offset  Size  Field
------  ----  ---------------------------------------------------------
0       6     TS         (timestamp_ms, big-endian)
6       1     VARIANT    (= 0x03 for PQHash v1)
7       2     COUNTER    (monotonic per generator, per ms)
9       64    TAIL       (SHAKE-256(content, 64 bytes) = 512 bits)
\end{lstlisting}

\textbf{Encodings:} 73 bytes encode to \textbf{98 characters} of canonical Base64URL \citep{rfc4648} (no padding) or \textbf{117 characters} of Crockford Base32 \citep{crockford}. Both encodings round-trip to identical 73-byte binary; binary sort and Crockford ASCII sort both reflect generation time. The Crockford form preserves case-insensitive transcription and accepts the standard \code{I/L $\to$ 1} and \code{O $\to$ 0} confusable foldings on decode; encode emits the canonical alphabet only.

\textbf{Security target:} the 512-bit SHAKE-256 tail provides $2^{256}$ classical second-preimage and birthday-collision resistance, dropping to $2^{256}$ Grover-second-preimage \citep{grover1996} and $\sim 2^{170}$ Brassard--H{\o}yer--Tapp \citep{bht1998} collision under a quantum adversary --- well above the $128$-bit post-quantum margin that motivates \pqid{}'s 256-bit random tail. The timestamp and counter are \emph{provenance metadata}, not security claims: two callers hashing the same content from different times produce identifiers that share the 64-byte tail and differ only in the 9-byte header.

\begin{calloutinfo}[At a glance]
\begin{itemize}
\item \textbf{Content-addressed, not random.} Two \pqhash{}es of the same content share the 64-byte tail; the 9-byte header gives a sortable, generator-aware view.
\item \textbf{Not a hash function.} \pqhash{} is a wire format \emph{built on} SHAKE-256; it is not itself a primitive you would substitute for SHA-3 in a protocol.
\item \textbf{Not a MAC, not a signature, not an envelope.} \pqhash{} provides no authentication and no confidentiality; pair with \pqcrypt{} (\cref{sec:monorepo}) when those properties are required.
\item \textbf{Inherits monorepo discipline.} Spec-first LaTeX; Python is the readable reference; C pins the canonical ABI; cross-language KAT vectors gate CI.
\end{itemize}
\end{calloutinfo}
